% =====================================================================
% Related work: Abhishek's paragraph (papers used in Stage 1)
% Paste into the shared Related Work section; the others add one paragraph each.
% =====================================================================
\paragraph{Pneumonia detection from chest X-rays.}
CheXNet~\cite{rajpurkar2017chexnet} fine-tuned an ImageNet-pretrained 121-layer DenseNet on the ChestX-ray14 dataset
of more than 100,000 frontal chest X-rays and reported pneumonia detection performance that exceeded the average of four practising radiologists.
It also visualised its predictions with class activation maps, an early step towards checking where such models look.
The paediatric dataset we use was introduced by Kermany \etal~\cite{kermany2018identifying}, who applied transfer learning from ImageNet to separate normal and pneumonia X-rays,
and it is widely used as a benchmark through its Kaggle release~\cite{mooney2018kaggle}.
Later work has questioned what such models actually learn: Zech \etal~\cite{zech2018variable} found that pneumonia classifiers exploit hospital-specific signals and lose accuracy on data from other hospitals,
and DeGrave \etal~\cite{degrave2021ai} found that COVID-19 classifiers rely on markers and regions outside the lungs, examples of shortcut learning in the sense of Geirhos \etal~\cite{geirhos2020shortcut}.
Our Stage~1 builds the baseline and the data pipeline for testing exactly this on the Kermany data. Unlike much previous work on this dataset, it uses a patient-level validation split and keeps the official test set untouched.
